% Generated draft: requires mathematical and source review.
\documentclass[11pt]{article}
\usepackage[margin=25mm]{geometry}
\usepackage{fontspec}
\setmainfont{Latin Modern Roman}
\usepackage{amsmath,amssymb,mathtools,unicode-math}
\setmathfont{Latin Modern Math}
\usepackage{enumitem,needspace,xurl,hyperref}
\setlength{\parindent}{0pt}
\setlength{\parskip}{5pt}
\setlength{\emergencystretch}{3em}
\newcommand{\N}{\mathbb{N}}
\newcommand{\Z}{\mathbb{Z}}
\newcommand{\Q}{\mathbb{Q}}
\newcommand{\R}{\mathbb{R}}
\newcommand{\C}{\mathbb{C}}
\newcommand{\F}{\mathbb{F}}
\newcommand{\A}{\mathbb{A}}
\newcommand{\PP}{\mathbb{P}}
\newcommand{\E}{\mathbb{E}}
\newcommand{\eps}{\varepsilon}
\newcommand{\dd}{\,\mathrm{d}}
\begin{document}
\section*{UnsolvedMath ID: 30001230}
\textbf{Source problem code:} catalogue entry
\subsection{Definitions and mathematical statement}
Let \(\Delta\subset\C\) be a disk centered at \(0\), and write \(X_t=\pi^{-1}(t)\). A compact complex manifold is Kähler if it admits a smooth, real, \(d\)-closed, positive definite \((1,1)\)-form. A compact complex manifold belongs to Fujiki class \(\mathcal C\) if it is bimeromorphic to a compact Kähler manifold.

\textbf{Deformation question in the supplied cleaned statement.} For every proper holomorphic submersion \(\pi:X\to\Delta\) with compact complex fibers, is the following implication true?
\[
\bigl( X_t\text{ is Kähler for every }t\in\Delta\setminus\{0\}\bigr)
\quad\Longrightarrow\quad X_0\in\mathcal C.
\]

\textbf{Current criterion attributed in the supplied record to Demailly's transcendental Morse conjecture.} Let \(Y\) be a compact complex manifold of complex dimension \(n\), and let \(T\) be a real, \(d\)-closed, positive \((1,1)\)-current on \(Y\). Positivity means that \(T\) acts nonnegatively on smooth strongly positive test forms of bidegree \((n-1,n-1)\). Write \(T_{\mathrm{ac}}\) for the absolutely continuous part of its measure coefficients with respect to a smooth positive volume form; this gives an almost everywhere defined positive semidefinite \((1,1)\)-form. The expression \((T_{\mathrm{ac}})^n\) denotes its pointwise top exterior power, integrated as a nonnegative measurable density. This definition is independent of the chosen smooth positive volume form.

A Kähler current on \(Y\) is a real, \(d\)-closed \((1,1)\)-current \(S\) for which there exist a smooth Hermitian \((1,1)\)-form \(\omega\) and a constant \(\varepsilon>0\) such that \(S\geq\varepsilon\omega\) in the sense of currents. Must
\[
\int_Y (T_{\mathrm{ac}})^n>0
\quad\Longrightarrow\quad
\exists\,S\text{ a Kähler current on }Y
\text{ with }[S]_{\mathrm{dR}}=[T]_{\mathrm{dR}}
\]
hold without assuming that \([T]_{\mathrm{dR}}\in H^2_{\mathrm{dR}}(Y,\R)\) is integral? Here integrality means membership in the image of \(H^2(Y,\Z)\to H^2_{\mathrm{dR}}(Y,\R)\). Equality of de Rham classes is understood using current representatives: \(S-T=dU\) for a real current \(U\) of degree one. No equality of Bott--Chern classes is added. The supplied condition does not explicitly require the displayed integral to be finite.

\textbf{Qualifications from the original wording.} The original statement describes a holomorphic family whose noncentral fibers are \emph{projective}, and asks, after obtaining a weak limit current \(T\) on \(X_0\), for Morse inequalities producing a Kähler current \(S\) in the same de Rham class as \(T\). It does not specify the approximating currents, their normalization, the hypotheses satisfied by their limit, or the quantitative inequalities requested. Consequently, that limiting-current step cannot be stated here as a fully specified additional assertion. The cleaned statement replaces projectivity by the broader Kähler hypothesis and presents the current criterion as a more general question. These formulations are retained as supplied; their equivalence and the claimed implication between the current criterion and the deformation question are not established by this draft.
\subsection{Short English statement}
The supplied cleaned statement asks whether a smooth proper holomorphic family with Kähler fibers away from the central parameter has central fiber in Fujiki class \(\mathcal C\). It also asks whether a closed positive \((1,1)\)-current with positive absolutely continuous top self-intersection has a Kähler current in its de Rham class, without an integrality assumption. The original wording uses projective noncentral fibers and an incompletely specified weak limit construction. This is a statement-only draft requiring human mathematical review.
\subsection{Sources}
\begin{itemize}
\item Supplied local UnsolvedMath registry record, permanent ID \(30001230\), title \emph{Kähler Limits and Transcendental Morse Inequalities}, exact source problem code \texttt{catalogue entry}, in the collection \emph{Oberwolfach Reports Open Problems}; proposed year: 2009. No individual proposer is supplied. The cleaned record attributes the transcendental Morse conjecture to Demailly.
\item Supplied report citation and literature source labeled ``original report'': \emph{Multiplier Ideal Sheaves in Algebraic and Complex Geometry} (2009), Oberwolfach Reports, pp. 1101--1156, DOI 10.4171/owr/2009/21. \url{https://doi.org/10.4171/owr/2009/21}. No exact problem page or passage within the report is supplied.
\item Supplied literature source labeled ``current status'': \url{https://arxiv.org/abs/2409.19957}. The record supplies no title, authors, or pinpoint citation for this link.
\item Supplied literature source labeled ``Popovici Ugarte'': \url{https://doi.org/10.1112/plms.12110}. The record supplies no article title or pinpoint citation for this link.
\item Supplied dated literature curation, checked exactly \texttt{2026-08-21}: status \texttt{partially\_solved}; assessment: the first limiting-current step is settled and the integral case is known, while the transcendental step and full deformation conjecture remain open. This assessment is supplied dated curation, not independently verified here. The source background explicitly requests independent verification of time-sensitive status and cited claims.
\item The registry labels the cleaned statement \texttt{corrected\_verified} and reports that it was checked against the official Oberwolfach report. This is a supplied source-verification claim, not an independent verification performed for this draft. The record was updated on \texttt{2026-09-16T12:06:53Z}.
\end{itemize}
\end{document}
